\section{Cardano's cubic formula}
\label{sec:cubic}

\subsection{Depressing the cubic}

Let
\[
 f(X)=aX^3+bX^2+cX+d,
 \qquad a\ne0.
\]
Put
\[
 A=\frac ba,\qquad B=\frac ca,\qquad C=\frac da,
\]
so that $f(X)=a(X^3+AX^2+BX+C)$.  Make the translation
\[
 X=Y-\frac A3.
\]
Direct expansion gives
\[
 X^3+AX^2+BX+C=Y^3+pY+q,
 \tag{C1}\label{eq:cubic-depress}
\]
where
\[
 p=B-\frac{A^2}{3}
   =\frac{3ac-b^2}{3a^2},
 \qquad
 q=\frac{2A^3}{27}-\frac{AB}{3}+C
   =\frac{2b^3-9abc+27a^2d}{27a^3}.
 \tag{C2}\label{eq:cubic-pq}
\]
The denominators are nonzero in characteristic zero.  It remains to factor the
depressed cubic
\[
 g(Y)=Y^3+pY+q.
\]

\subsection{The compatible radical equations}

Choose $u,v$ satisfying
\[
 u^3+v^3=-q,
 \qquad
 uv=-\frac p3.
 \tag{C3}\label{eq:cardano-compatible}
\]
Then
\[
 (u+v)^3+p(u+v)+q
 =u^3+v^3+3uv(u+v)+p(u+v)+q=0.
\]
Thus $u+v$ is a root.  The product equation in
\eqref{eq:cardano-compatible} is essential.  Independent choices of cube-root
branches do not in general satisfy it.

Let $\omega$ be a primitive cube-root witness in the algebraic form
\[
 \omega^2+\omega+1=0.
 \tag{C4}\label{eq:omega}
\]
Then $\omega^3=1$ and $1+\omega+\omega^2=0$.  In characteristic zero,
$\omega\ne1$.  Such a witness itself uses only a square radical: if
$\tau^2=-3$, then $\omega=(-1+\tau)/2$ satisfies~\eqref{eq:omega}.

Define
\[
 r_0=u+v,
 \quad
 r_1=\omega u+\omega^2v,
 \quad
 r_2=\omega^2u+\omega v.
 \tag{C5}\label{eq:cardano-roots-depressed}
\]

\begin{theorem}[Cardano factorization]
\label{thm:cardano-factorization}
From~\eqref{eq:cardano-compatible} and~\eqref{eq:omega} one can derive
\[
 Y^3+pY+q=(Y-r_0)(Y-r_1)(Y-r_2).
 \tag{C6}\label{eq:cardano-factorization}
\]
Consequently the roots of the original cubic are
\[
 x_k=r_k-\frac A3,
 \qquad k=0,1,2.
 \tag{C7}\label{eq:cardano-original-roots}
\]
\end{theorem}

\begin{proof}
The sum of the three displayed values is
\[
 r_0+r_1+r_2
 =(1+\omega+\omega^2)u+(1+\omega+\omega^2)v=0.
\]
A direct expansion, reduced by $\omega^2+\omega+1=0$, gives
\[
 r_0r_1+r_0r_2+r_1r_2=-3uv=p.
\]
Similarly,
\[
 r_0r_1r_2=u^3+v^3=-q.
\]
The monic polynomial with these three roots is therefore
\[
 Y^3-(r_0+r_1+r_2)Y^2
 +(r_0r_1+r_0r_2+r_1r_2)Y-r_0r_1r_2
 =Y^3+pY+q.
\]
Combining this identity with~\eqref{eq:cubic-depress} proves the claim.
\end{proof}

The factorization, rather than three independent substitution calculations,
is the preferred certificate.  It proves all roots simultaneously and retains
multiplicity when two or three values coincide.

\subsection{Constructing a compatible pair without a hidden pole}

The usual displayed formula introduces
\[
 \Delta=\left(\frac q2\right)^2+\left(\frac p3\right)^3
 \tag{C8}\label{eq:cardano-delta}
\]
and a square-root witness $s^2=\Delta$.  Put
\[
 R=-\frac q2+s,
 \qquad
 S=-\frac q2-s.
 \tag{C9}\label{eq:RS}
\]
Then
\[
 R+S=-q,
 \qquad
 RS=\frac{q^2}{4}-s^2=-\left(\frac p3\right)^3.
 \tag{C10}\label{eq:RS-relations}
\]
A textbook often says ``take $u^3=R$ and $v^3=S$'' without controlling the
branches.  The following finite construction is correct in every case.

\paragraph{Case $p\ne0$.}
Equation~\eqref{eq:RS-relations} implies $R\ne0$ and $S\ne0$.  Adjoin a cube
root $u$ with $u^3=R$ and define
\[
 v=-\frac{p}{3u}.
 \tag{C11}\label{eq:v-defined}
\]
The denominator is nonzero.  Then $uv=-p/3$, and
\[
 v^3=\frac{(-p/3)^3}{u^3}
 =\frac{-(p/3)^3}{R}=S
\]
by $RS=-(p/3)^3$.  Hence $u^3+v^3=R+S=-q$.

\paragraph{Case $p=0$.}
Now $RS=0$.  Adjoin cube-root witnesses independently,
\[
 u^3=R,\qquad v^3=S.
\]
Then
\[
 (uv)^3=u^3v^3=RS=0.
\]
A field has no nonzero nilpotents, so $uv=0=-p/3$; also
$u^3+v^3=R+S=-q$.  No decision whether $R$ or $S$ vanishes is needed.

Thus the construction uses one square radical, at most two cube radicals, and
the square radical used to construct $\omega$.  The $p\ne0$ branch needs only
one independently chosen cube root because the other is defined by division;
the $p=0$ branch uses two cube-root witnesses but contains no division by a
possibly zero quantity.

\subsection{The familiar displayed formula}

With a compatible choice of cube roots, the three roots may be written
\[
 x_k=-\frac{b}{3a}
 +\omega^k\sqrt[3]{-\frac q2+\sqrt\Delta}
 +\omega^{-k}\sqrt[3]{-\frac q2-\sqrt\Delta},
 \qquad k=0,1,2,
 \tag{C12}\label{eq:classical-cardano}
\]
where the two cube radicals are constrained so that their product is $-p/3$.
Equation~\eqref{eq:v-defined} is a branch-safe implementation of that
constraint when $p\ne0$; the $p=0$ branch above removes the apparent pole.

\subsection{Constructive exhaustiveness}

If $g(y)=0$, equation~\eqref{eq:cardano-factorization} gives
\[
 (y-r_0)(y-r_1)(y-r_2)=0.
\]
In a coded field, test $y-r_0=0$.  If false, cancel its explicit inverse and
test $y-r_1=0$; if that is false, cancellation gives $y-r_2=0$.  This is a
finite decision tree, not an application of excluded middle to an arbitrary
proposition.  The same argument after translation proves that every root of
$f$ occurs in~\eqref{eq:cardano-original-roots}.

\subsection{Primitive-recursive verification}

The certificate checker expands the right-hand side of
\eqref{eq:cardano-factorization} and reduces it by the three relations
\[
 \omega^2+\omega+1=0,
 \qquad
 3uv+p=0,
 \qquad
 u^3+v^3+q=0.
\]
Because the degrees and number of variables are fixed, even a naive sparse
normalizer is primitive recursive.  The supplied verification script computes
a Groebner normal form only as an independent convenience check; the formal
iPRA proof needs merely a fixed sequence of ring rewrites, exactly as the
existing Lean theorem \code{cardano_factorization} does.  ProveIt's
\code{cubic_eq_zero_iff} already packages the semantic soundness and
exhaustiveness statement over a characteristic-zero field.
